%%KEEP 1-87%%
\section{Introduction}
\label{sec:introduction}

%%KEEP 91-98%%

Runtime-series benchmarking spans classical statistical, neural and zero-shot
foundation models \citep{DiMenna2026RuntimeForecasting}, and cloud studies
combine multivariate resource inputs, recurrent architectures and joint
multi-output prediction \citep{Ullah2023CloudBiLSTM,Tripathi2025MMRP}.
\emph{Contextual input} supplies other resources' histories to a target model
without requiring shared parameters; \emph{representation sharing} learns
several targets jointly. Multi-task learning offers inductive transfer through
related tasks \citep{Caruana1997Multitask}, but task combinations and sharing
arrangements require evaluation
\citep{Standley2020TaskGrouping,Fifty2021TaskGroupings,Nochumsohn2025MTLinear}.
Our bounded contribution is a within-study comparison that separates the two
decisions, not a new forecasting architecture (Section~\ref{sec:related-work}).

%%KEEP 123-183%%

\section{Related Work}
\label{sec:related-work}

Runtime telemetry supports configuration and resource-management decisions in
cloud and self-adaptive software systems
\citep{Taherizadeh2018Monitoring,Bauer2020SelfAware}, and prediction-based
provisioning and service-quality forecasting connect forecast accuracy with
software design \citep{Calheiros2015ARIMA,Syu2017DynamicQoS}. Forecast quality
nevertheless differs from controller effectiveness \citep{Qu2018Autoscaling};
the question here concerns input selection and model management, not a
measured autoscaling benefit. Multivariate covariates, hybrid resource
predictors, joint outputs and multi-provider evaluation are established
\citep{Tran2018MultivariateFuzzy,Xu2022esDNN,Tripathi2025MMRP,
Rossi2025CloudUncertainty,Zhang2025SWTCLSTM}, alongside online resource and
microservice-indicator forecasting \citep{Gao2026OMCR,Gao2026MaMiClif}.
\citet{DiMenna2026RuntimeForecasting} fit conventional models independently
per series and apply foundation models zero-shot, whereas
\citet{Ullah2023CloudBiLSTM} forecast ten channels jointly with a common
BiLSTM on fastStorage and Rnd. TimeTrack's reported resource correlations
\citep{Meliani2025TimeTrack} motivate contextual inputs but do not establish
predictive value conditional on the target's own history.

Input channels, output targets and parameter sharing are separate choices.
MMRP compares univariate and multivariate single-target forecasts
\citep{Tripathi2025MMRP}, a substantive S1--S2-related precedent without a
joint resource-sharing arm, and Ullah's joint BiLSTM shares a representation
without an independent all-input target-model bridge. PatchTST reuses
embedding and Transformer weights across independently processed channels
\citep{Nie2023PatchTST}: channel independence is shared parameters with
own-channel input, the arrangement our design calls M0. DLinear forecasts
trend and residual components with linear maps \citep{Zeng2023DLinear}, and
iTransformer models relationships between variate tokens through attention
\citep{Liu2024iTransformer}; tier~2 evaluates declared adaptations of all
three. DeepAR and global forecasting share functions across series rather
than presenting them as simultaneous covariates
\citep{Salinas2020DeepAR,MonteroManso2021GlobalForecasting}.

%%KEEP 260-272%%

\citet{Nikentari2022NARMAXLSTM} combine lagged station inputs with shared LSTM
layers and discuss negative transfer, MTLinear analyses gradient conflict,
balance and grouping for variate groups \citep{Nochumsohn2025MTLinear}, and
\citet{Bi2019CloudMultitask} report training-speed comparisons for cloud
multi-task forecasting. None establishes an independent all-input
target-model bridge or a measured target-consolidation trade-off;
Supplementary Table~[[SUPP:tab:related-work-comparison]] codes these designs
against ours (Supplementary Section~S-A).

The contribution is therefore the within-study comparison, not an
architecture or priority claim. S2 is the intermediate control: each target
model receives all resource histories but keeps separate parameters, so S1--S2
examines contextual input and S2--M1 the move to a shared predictor.
Capacity, optimisation and selected settings also differ, so these are not
pure causal interventions. Tier~2 adds M0 to complete the $2\times2$ on a
reduced Bitbrains cohort and repeats it across backbones, while post-hoc
common-context, capacity and Ridge controls qualify the tier-1 contrasts
(Section~\ref{sec:posthoc-controls}).

%%KEEP 414-468%%

Models are entity-specific: for each machine or VM, outer fold and seed, S1/S2
fit four target-specific networks and M1 (and, in tier~2, M0) one joint
four-resource network, never shared across entities. Contexts contain
observations available at or before $t$; targets lie strictly afterwards, and
matched origins, targets, horizons and partitions prevent sample differences
from changing a contrast. The four resources are CPU utilisation, aggregate
throughput, memory usage and disk-write activity, and every formulation
predicts its full horizon vector directly.

\subsection{Four Formulations}
\label{sec:formulations}
\label{sec:s1}
\label{sec:s2}
\label{sec:m1}
\label{sec:m0}

S1 gives each target resource $r$ its own model on its own history,
%%KEEP 508-521%%
S2 gives each target its own model on all four histories,
%%KEEP 532-540%%
so non-target histories act as covariates, but gradients from one target
cannot update another target's model. M1 combines all-resource input with an
encoder shared across targets and resource-specific heads:
%%KEEP 559-568%%
%%KEEP 576-592%%
Equal coefficients do not imply equal realised losses or gradients, and no
task weights are learned. M0 applies one shared forecaster to each resource's
own history,
%%KEEP 607-614%%
trained with the same objective; it is evaluated only in tier~2, with
backbone-specific forms in Section~\ref{sec:tier2-design}. Supplementary
Section~S-B0 gives the full formulation text.

%%KEEP 624-649%%
None presumes an error ordering, and native-unit errors are never averaged
across resources. On matched support the S1-to-M1 raw difference equals the
sum of its two intermediate differences, without implying mechanistic
additivity; percentages are not additive because each uses its own reference.

%%KEEP 659-728%%

\subsection{Data, Targets and Leakage Control}
\label{sec:datasets}
\label{sec:preprocessing}
\label{sec:leakage-control}
\label{sec:population-restrictions}
\label{sec:replication-design}

Table~\ref{tab:dataset-summary} summarises the sources. TimeTrack, telemetry
from the OpenAirInterface (OAI) CI/CD system on seven machines between
24 June and 19 July 2024, is framework-development evidence with a measured
median sampling interval of 42.285164 seconds; gaps longer than 84.57 seconds
start a new segment. The Bitbrains fastStorage and Rnd traces are two separate
external environments with five-minute sampling, in which physically invalid
rows are removed and gaps longer than 450 seconds delimit segments. Materna
T1--T3 are consecutive, dependent blocks from one source on a 300-second grid
and provide secondary temporal replication, not independent datasets. All
sources map to CPU, throughput, memory and disk write: CPU and memory use the
identity transform and throughput and disk write use $\log(1+x)$; each
transformed series is robust-scaled with training-only median and
interquartile range, which must be positive (no epsilon or fallback), and is
returned to native units with a non-negativity projection for throughput and
disk write.

External eligibility requires structural validity and a positive
transformed-resource interquartile range in every required training
partition. From each eligible pool, 35 VMs were selected by ascending SHA-256
rank of a frozen key, without outcome information; for Materna, 74 VMs passed
the joint rule and the first 35 are used in all three blocks. The cohorts
therefore represent scale-estimable, structurally eligible VMs, not every
archived VM.

An origin is legal only when the maximum context and every future target lie
in one contiguous segment, and common support is built at the largest
context, so all formulations share origins, targets, folds, horizons and
seeds. Five expanding outer folds use the first half of the ordered legal
origins before the first of five successive ten-percent test blocks, and
training origins whose maximum-horizon target reaches the test block are
purged. The latest 15\% of outer-training origins form inner validation,
purged again; scalers, contexts and epochs use only training and validation
data, and each final model is refitted once on the complete outer-training
support and evaluated once on its test block. Horizons are
$\{1,2,4,8,16,32,64,128\}$ samples in TimeTrack (about 0.705 to 90.208
minutes) and $\{1,2,4,9,18\}$ samples externally (5, 10, 20, 45 and 90
minutes), with candidate contexts of 32, 64 and 128 and of 4, 9 and 18
samples. Bitbrains environments and Materna blocks are analysed separately
before any equal-weight descriptive summary, and no score pools sources.
Supplementary Section~S-B1 gives the full data, preprocessing and leakage
protocol.

%%KEEP 795-844%%

\subsection{Models, Training and Selection}
\label{sec:architecture-training}
\label{sec:model-selection}

Tier~1 uses a plain two-layer unidirectional LSTM \citep{Hochreiter1997LSTM}
with hidden size 64 and dropout 0.1. S1 and S2 minimise target-specific scaled
mean squared error and M1 the mean of the four resource-specific errors
(Equation~\ref{eq:m1-loss}). Training uses Adam with learning rate 0.001,
batch size 256, at most 50 epochs, early-stopping patience 7, gradient-norm
clipping at 1.0 and seeds 0, 1 and 2, with no architecture or optimiser
search. Context length is selected from the candidate grid on
inner-validation scaled MAE at each seed's MSE-selected epoch, per target for
S1 and S2 and per entity and fold for M1, and exact ties favour the shorter
context; each final model is refitted with the selected context and
seed-specific epoch count. In the five-horizon external design the
complements have 203,028 (S1), 206,100 (S2) and 52,500 (M1) parameters.
Reported fit time sums successful worker fit-hours and is not elapsed wall
time, energy or deployment cost (Supplementary Section~S-B2).

\subsection{Evaluation and Uncertainty}
\label{sec:evaluation}

Primary evaluation uses native-unit MAE after inverse transformation and
projection; native errors are never averaged across different resources, and
RMSE, $R^2$ and native NMAE are secondary diagnostics. Persistence supplies a
separate absolute forecasting check:
%%KEEP 1038-1045%%
averaged over entity--fold cells and then horizons, with undefined
zero-denominator skills excluded and counted. MAE is averaged over the three
seeds within each matched entity, fold, resource, horizon and formulation
cell; for a fixed resource and horizon, seed-averaged cells receive equal
weight to form MacroMAE, and a resource-overall MacroMAE gives equal weight to
every horizon. For candidate $C$ and baseline $B$, the reported gain is
%%KEEP 1061-1075%%

Prospectively frozen classes---\textsc{Broad Positive},
\textsc{Conditionally Positive}, \textsc{Broad Negative} and
\textsc{Neutral or Mixed}---combine the overall gain (at least $+1\%$,
strictly positive, or at most $-1\%$) with consistency over horizons,
entities, folds and seeds; the thresholds are in Supplementary
Section~S-B3 and the supplementary decision criteria. Intervals never enter
the classes, and no threshold is adapted after observing outcomes. External
uncertainty uses a paired VM-cluster bootstrap within each Bitbrains
environment and Materna trace that samples 35 VMs with replacement, retains
all five folds and the paired formulation structure, and uses 10,000
replicates, percentile 95\% intervals and random seed 20261007. TimeTrack has
no such intervals.

\subsection{Post-Hoc Sensitivity Analyses}
\label{sec:posthoc-controls}

Three analyses were added after the original experiment was frozen. They are
post-hoc sensitivity controls, not amendments to the prospective design or
replacements for its results, and retain the frozen cohorts, folds, horizons,
purging, support, transforms, scalers and native-unit evaluation.
\emph{Common-context S2} imposes the frozen H64-M1 context on every S2 target,
removing S2's context-selection freedom. \emph{H129 M1} raises the shared
hidden size from H64 to H129 (206,420 parameters externally), approximately
matching the 206,100-parameter S2 complement, with selection repeated afresh.
\emph{Matched-context Ridge AR/ARX} compares own-lag and all-resource-lag
Ridge models at the H64-M1 context, with $\alpha$ selected on inner
native-unit validation MAE. Neither LSTM control equalises architecture or
optimisation, the two are separate rather than a jointly matched
intervention, and Ridge tests input scope in a second model family rather
than isolating model family (procedures in Supplementary Section~S-B4).

%%KEEP 1203-1215%%

%%KEEP 1217-1229%%

\paragraph{Backbones.} Hyperparameters are fixed in advance and identical
across cells except where a cell's definition requires otherwise (learning
rate $10^{-3}$, batch size 256, at most 50 inner epochs, patience 7, clipping
at 1.0, seeds 0--2), with no hyperparameter search. The LSTM is as in tier~1,
with M0 an input-size-one encoder shared across the four own-resource
histories. DLinear \citep{Zeng2023DLinear} forecasts trend and residual
components with linear maps: its M0 reuses one pair of maps on every
resource channel without resource-specific outputs, S2 is a decomposed linear
ARX, and M1 passes all inputs through a shared linear bottleneck of width 16
before per-target heads. The PatchTST-style backbone \citep{Nie2023PatchTST}
applies one patch Transformer to each channel in M0 and channel-dependent
tokens in S2 and M1. The iTransformer-style backbone
\citep{Liu2024iTransformer} embeds each resource history as a variate token;
its S1 and M0 cells are undefined because a single token reduces attention to
a token-wise multilayer perceptron. Both transformer-style backbones are
declared adaptations, not canonical implementations, and every new cell
passed structural checks before any data fit. Full specifications are in
Supplementary Section~S-B5.

%%KEEP 1275-1314%%
D-1 is the maximum absolute lagged cross-correlation (lags 1--18) between
ordered resource pairs on inner-training series; D-2 reuses the post-hoc Ridge
ARX:AR gain; D-3 is the Spearman correlation across the ten VMs of each block
between D-1 or D-2 and the LSTM entity log effects
$D(C\!:\!B)=\ln A_C-\ln A_B$; D-4 measures, in the LSTM M1 shared encoder at
the selected epoch over eight keyed batch pairs, the pairwise task-gradient
cosines, the ratio of the largest to the smallest task-gradient norm and the
excess sign-conflict rate $K$; D-5 fits all-input LSTM models on pairwise and
leave-one-out resource subsets, from which D-6 derives directed transfer
$F_i(j)$ and leave-one-out contribution $Q_{i,j}$ (positive values favour
adding or retaining resource $j$); and D-7 is each task's share of the M1
training loss. Definitions are in Supplementary Section~S-B6.
%%KEEP 1356-1360%%

\section{Results}

Figure~\ref{fig:primary-gains} reports the frozen tier-1 resource-overall
MacroMAE gains for each evidence block; positive values favour the candidate.
Consistency classes and VM-cluster intervals are reported separately, and the
post-hoc controls supplement rather than replace them (Supplementary
Tables~S1--S2 and Table~\ref{tab:posthoc-controls}). Tier-2 results come from the
separately frozen extension of Section~\ref{sec:tier2-design}; they are
descriptive, never replace tier-1 results and are reported in separate
paragraphs.

%%KEEP 1393-1474%%
RQ2 first compares S2 with S1. Figure~\ref{fig:primary-gains}(a) shows
selective contextual value, not uniform improvement. In tier~1, memory had the
clearest recurring positive pattern: it was broadly positive in fastStorage
and Materna T2, conditionally positive in Rnd and T3 and the sole positive
resource in TimeTrack, but approximately neutral and slightly negative in T1
($-0.02\%$). The other resources changed direction across blocks, Rnd
throughput and T1 disk write being adverse where other blocks were near zero
or positive, and many intervals crossed zero. The matched-context Ridge
ARX:AR comparison also showed selective, family-dependent value: memory
improved under Ridge in all five external blocks but reversed in TimeTrack,
and the Ridge TimeTrack throughput gain of $-838.45\%$ is concentrated in one
machine/fold ($99.93\%$ of positive excess error) rather than a broad
failure. Existing-seed diagnostics find sign changes in 9/24 LSTM S2:S1
block/resource summaries (Supplementary Table~S2), so seed-varying point gains
are not seed-invariant improvements. Block-by-block classes, intervals and the
Ridge audit are in Supplementary Section~S-C1.

%%KEEP 1526-1548%%
RQ3 first compares M1 with S2. Throughput, memory and disk write had negative
primary gains in all six blocks (Figure~\ref{fig:primary-gains}(b)); CPU was
positive only in TimeTrack, fastStorage and Materna T1, and these small
exceptions do not support uniformly beneficial sharing. Rnd and Materna T2
were broadly negative for every resource with intervals below zero, and most
other throughput, memory and disk-write cells were broadly negative. Under
both post-hoc controls, H64 M1 against common-context S2 and H129 M1 against
S2, throughput, memory and disk-write point estimates remained adverse in all
six blocks (Table~\ref{tab:posthoc-controls}), so
context-selection freedom and the original capacity difference are
incomplete explanations, although the controls do not identify sharing alone
as the cause. In all 24 original M1:S2 block/resource summaries the
matched-cell median MAE ratio exceeded one and the win rate was below $50\%$
(Supplementary Table~S1), yet magnitudes were tail-concentrated: Materna T1/T3
disk-write median ratios were approximately $1.0246/1.0347$ despite MacroMAE
gains of $-27.46\%/-71.30\%$, with VM266 and VM272 contributing $82.10\%$ and
$86.97\%$ of positive excess error. Existing-seed diagnostics find sign
changes in 4/24 M1:S2 summaries, and sign stability need not imply magnitude
stability (Supplementary Section~S-C2).

%%KEEP 1621-1732%%
\subsection{RQ3 Complete Joint-System Change}
\label{sec:results-rq3}

The complete joint change, M1 versus S1 (Figure~\ref{fig:primary-gains}(c)),
was negative in every displayed cell except Materna T2 CPU ($+1.75\%$, with an
interval crossing zero), and Rnd was broadly negative for every resource. Its
gains are reported directly, not added from the S2:S1 and M1:S2 percentages
(Section~\ref{sec:decomposition}); classes and intervals are in Supplementary
Section~S-C4.

%%KEEP 1762-1765%%
penalty (Figure~\ref{fig:tier2-mechanism}, Supplementary
Table~[[SUPP:tab:tier2-diagnostics]]). None was used to change any model, and
none establishes a cause.

%%KEEP 1776-1828%%
\subsection{RQ5 Resource, Environment, and Temporal Heterogeneity}
\label{sec:results-rq4}
\label{sec:cross-dataset-synthesis}

Heterogeneity appears at each prespecified level
(Figure~\ref{fig:primary-gains}). For S2:S1, CPU, throughput and disk write
change direction across evidence blocks, whereas memory is positive in five
and nearly zero and negative only in Materna T1. For sharing, fastStorage CPU
was near zero while the other three resources were negative, Rnd was broadly
negative for all four, and Materna M1-versus-S2 CPU changed from positive in
T1 to negative in T2--T3. The prespecified original-RQ4 criterion for temporal
heterogeneity was met: seven of the 12 resource--comparison pairs changed
frozen classification across T1--T3 and five changed raw sign. This is a
descriptive criterion, not an inferential claim (Supplementary
Section~S-C4).

Table~\ref{tab:cross-study-synthesis} retains the frozen TimeTrack,
equal-environment Bitbrains and equal-temporal-block Materna summaries without
a global score, with the tier-2 counts in a separate row. External memory has
favourable LSTM/Ridge contextual agreement, whereas throughput reverses
direction between families, and Materna's consecutive blocks share VM
identifiers, so their agreement is not independent replication.

%%KEEP 2025-2077%%
\subsection{Persistence Boundary}
\label{sec:results-persistence}

Relative formulation gains do not establish superiority to a persistence
forecast. In TimeTrack, S1, S2 and M1 had positive mean skill for CPU,
throughput and disk write, but memory was the boundary ($-0.1272$, $-0.1494$
and $-0.6047$), so the large M1-versus-S2 memory penalty coincided with weak
absolute performance. Overall CPU skill was negative for S1, S2 and M1 in both
Bitbrains environments, and neither environment contained a formulation that
beat persistence for every resource. In Materna, S2 beat persistence in all 12
trace--resource overall aggregates, whereas M1 failed the check for T1 and T3
throughput and disk write.
%%KEEP 1951-1955%%

\subsection{Accuracy-Efficiency Trade-Off}
\label{sec:results-efficiency}

The original H64 M1 consolidates the target-model complement for one entity,
outer fold and seed (Table~\ref{tab:efficiency}): one joint network and
52,500 parameters against four networks and 206,100 parameters for S2,
approximately $74.5\%$ fewer parameters and $75\%$ fewer models. The frozen
Materna totals are 85.594240 worker fit-hours for S1, 75.303186 for S2 and
20.405083 for M1; they are not elapsed wall time, energy or deployment cost,
and the post-hoc H129 comparison does not inherit H64's parameter-reduction
claim. In tier~2, shared LSTM, PatchTST-style and iTransformer-style models
used roughly a quarter of their independent complements' parameters, whereas
DLinear M1 was not smaller than DLinear S2 at contexts 4 and 9
(Section~\ref{sec:results-rq2}). These local inventories coexist with the
accuracy penalties above and establish no production savings.

%%KEEP 1979-2003%%
\section{Discussion}
\label{sec:discussion}

\subsection{When Cross-Resource Context Helps}
\label{sec:discussion-context}

Cross-family agreement for external memory, TimeTrack's reversal and T1's
seed-variable estimate support selective contextual value, not a universal
memory benefit, in line with MMRP's non-uniform contextual gains
\citep{Tripathi2025MMRP} and the relevance of environment-specific evaluation
\citep{Rossi2025CloudUncertainty}. TimeTrack correlations
\citep{Meliani2025TimeTrack} motivate inputs, not conditional predictive value
or causal resource coupling.
%%KEEP 2143-2152%%
The separate sensitivity controls retain adverse throughput, memory and
disk-write points across all six blocks, and direction need not imply uniform
severity: Materna T1/T3 disk-write median ratios are about $1.025/1.035$,
despite much larger tail-sensitive MacroMAE penalties.

%%KEEP 2163-2199%%
\subsection{Implications and the Accuracy-Consolidation Trade-Off}
\label{sec:discussion-systems}
\label{sec:discussion-efficiency}

Forecasting supports adaptation
\citep{Taherizadeh2018Monitoring,Bauer2020SelfAware}, but adding telemetry and
consolidating predictors require separate resource- and horizon-specific
validation on future-only support, including persistence. Controller outcomes
depend on policy and platform \citep{Qu2018Autoscaling}; no autoscaler,
scheduler or controller was deployed, and production latency, energy,
monetary savings and SLA improvements were not measured. Local model and
parameter reductions coexist with accuracy losses and are backbone-specific,
so compactness is a trade-off, not a free predictive gain; joint
accuracy/computation assessment follows task-grouping and service-forecasting
research \citep{Standley2020TaskGrouping,Syu2017DynamicQoS}. Per-entity
consolidation may still be attractive where modest target-specific accuracy
penalties are acceptable, provided adverse tails are checked.

%%KEEP 2232-2268%%
\subsection{Construct Validity}
\label{sec:threats-construct}

Tier~1 has three of the four $2\times2$ cells; M0 exists only in tier~2, on a
reduced cohort, and tier-2 cells are defined by backbone structure:
iTransformer S1 and M0 are undefined, DLinear M0 ties one map across all four
resources without resource-specific outputs, and DLinear M1's width-16
bottleneck has lower rank than four S2 maps at context 4, so DLinear M1:S2
partly reflects reduced capacity. The transformer-style backbones are declared
adaptations. Original capacity and context differences are addressed
separately, not jointly, by the post-hoc controls, and resource mappings are
source-specific rather than perfectly equivalent. MacroMAE gains describe one
error estimand; metric choice can alter rankings \citep{Hyndman2006Accuracy},
and scaled-space training differs from native-unit evaluation.

\subsection{Internal Validity}
\label{sec:threats-internal}

Chronology, purging and training-only preprocessing constrain leakage
\citep{Cerqueira2020Evaluation,Hyndman2021Forecasting}, but matched support and
provenance cannot exclude all implementation or data errors. Original
external criteria were prospective; the later controls are post-hoc, and the
tier-2 protocol was frozen by a Git tag, package digest and execution seal
before any tier-2 run, not by a time-stamped external preregistration. Only
three neural seeds were recorded: 9/24 S2:S1 and 4/24 M1:S2 tier-1
block/resource summaries change sign across them (Supplementary Table~S2).
Tier~2 uses fixed, untuned hyperparameters and LSTM-selected H64-M1 contexts
for every backbone, so the modern backbones may be under-tuned; the largest
tier-2 aggregate was stable across seeds, with per-VM seed spreads outside
VMs 233 and 269 of at most 0.046 native units. Gradient diagnostics cover
only the LSTM at the selected epoch, D-2 reuses Ridge outcomes already
observed, and Ridge uses a finite five-alpha grid at LSTM-selected contexts.

\subsection{External Validity}
\label{sec:threats-external}

The sources represent particular runtime and cloud contexts, not probability
samples of all VMs, and Materna's blocks are repeated-within-source evidence.
Tier~2 covers 20 of the 70 Bitbrains VMs, with no Materna or TimeTrack data;
tier~1 remains the full-cohort evidence. Canonical or tuned transformer
forecasters, mixture-of-experts, selective grouping, adaptive sharing,
gradient surgery, loss balancing and foundation forecasters were not tested,
and four resource semantics do not establish accuracy for application
latency, QoS or downstream SLA outcomes.

\subsection{Conclusion Validity}
\label{sec:threats-conclusion}

%%KEEP 2360-2370%%
Paired VM-cluster bootstrap intervals condition on the evaluated cohorts and
are not classical significance claims; overlapping folds, multiple horizons,
repeated seeds and dependent Materna blocks are not IID replicates, and no
global effect or meta-analysis is used to strengthen these conclusions.
Supplementary Section~S-D keeps the full-length discussion and threats to
validity.

\section{Conclusion}
\label{sec:conclusion}

%%KEEP 2386-2391%%
Contextual input gave small and selective benefits: memory in the external
environments showed the most consistent favourable pattern, but directions
depended on the resource, the environment and the model family, and no
tier-2 context effect held across backbones.

Representation sharing carried a recurring accuracy cost. Tier-1 throughput,
memory and disk-write penalties remained in all six blocks under separate
common-context and approximately capacity-matched controls, and in tier~2
all-input sharing lowered accuracy in 29 of 32 backbone, resource and
environment cells, with own-resource sharing adverse throughout Rnd.
Magnitudes were not uniform: matched-cell median ratios were close to one in
many comparisons, and the largest tier-2 penalty, $-139.61\%$ for DLinear Rnd
CPU own-input sharing, is valid under the frozen aggregation but is dominated
by one high-scale VM while its negative direction is broad.

%%KEEP 2411-2420%%
Contextual input and model consolidation should therefore be evaluated
separately for each resource and forecasting design, checking typical-cell
behaviour, adverse tails and persistence skill within the eligible
population. The evidence does not support a general rejection of shared or
multi-task forecasting.

\section*{Data availability}

TimeTrack is publicly available on
\href{https://www.kaggle.com/datasets/abdelghanimeliani/open-air-interface-cicd-timeseries-dataset}{Kaggle}
(Apache 2.0 in its dataset metadata) and in the publisher's
\href{https://doi.org/10.5281/zenodo.15442251}{Zenodo archive
(doi:10.5281/zenodo.15442251; CC BY 4.0)}; its description paper is
\citet{Meliani2025TimeTrack}. Bitbrains fastStorage and Rnd are publicly
available as \href{https://atlarge-research.com/gwa-t-12/}{GWA-T-12}
\citep{Shen2015Bitbrains}, and Materna T1--T3 as
\href{https://atlarge-research.com/gwa-t-13/}{GWA-T-13}
\citep{Kohne2014FederatedCloudSim,Kohne2016SLAScheduling}. We acknowledge
Bitbrains IT Services Inc. and Materna GmbH Information \& Communications for
providing these traces through the Grid Workloads Archive; their source pages
specify acknowledgement and citation conditions rather than a standard
open-data licence, and readers should obtain the raw data from these sources.

Source code, protocols, derived analysis outputs and SHA-256 inventories for
the original experiments, post-hoc controls and robustness diagnostics are at
the
\href{https://github.com/KADOUMAAbdelhak/TimeTrack-Forecasting-Framework/tree/0e8cad884896a706f62f0fae4ab2b52812f9b033}{project repository checkpoint}
\texttt{\seqsplit{0e8cad884896a706f62f0fae4ab2b52812f9b033}}; fitted-model and
per-fit prediction collections are identified by checksum inventories but not
included. The tier-2 protocol, code, configuration, input digests, execution
seal and launcher are frozen by the annotated tag
\href{https://github.com/KADOUMAAbdelhak/TimeTrack-Forecasting-Framework/tree/fgcs-v04-min-final-r2}{\texttt{fgcs-v04-min-final-r2}}
(commit \texttt{\seqsplit{7ef468c0cc21c1c8ca33bfd9db0ee919e0dee24a}}), and
sealed inputs not held in Git are archived on Zenodo
(\href{https://doi.org/10.5281/zenodo.24189108}{doi:\allowbreak 10.5281/\allowbreak zenodo.24189108}).
The tier-2 analysis output (SHA-256
\texttt{\seqsplit{308841865ec87b04e10ecb64f20bc13ce2e671511a1574056028f7b5710dd17a}})
and the DLinear Rnd CPU audit summary (SHA-256
\texttt{\seqsplit{88b6f56a81d2db8eb87a4d878569c95595c247f1368dd2c8460429dbbb9a1b58}}),
with the closed tier-2 journal records and a file-level SHA-256 inventory, are
archived on Zenodo
(\href{https://doi.org/10.5281/zenodo.24339239}{doi:\allowbreak 10.5281/\allowbreak zenodo.24339239},
version v0.4-outputs; launcher records with local paths redacted, model
weights not included), and the tier-2 figures and tables are generated by
scripts that verify these digests. An earlier preregistration of a superseded
protocol version (\href{https://osf.io/kuwra}{OSF kuwra}) is retained as
historical evidence.

%%KEEP 2482-2518%%
